% ============================================================
% Chapter: Design Basis and Metocean
% Assessed MLOs: MLO-1, MLO-2
% Excellent-grade rubric criteria:
% - Design basis and philosophy quality (5%): the design philosophy is explicit,
%   technically justified, linked to failure modes, and translated into measurable
%   requirements and governing load cases.
% - Meteocean report quality (5%): the environmental basis is site-specific,
%   source-traceable, complete, and converted into defensible design conditions.
% ============================================================

\section{Design Basis and Metocean}
\pagecheck{6--7}

% Rubric Checklist: Design basis and philosophy quality (5%); Meteocean report quality (5%).
\instructionplaceholder{Project description and location}{Open the report with a concise project overview. State the location, infrastructure function, users served, and why a floating or submerged solution is being considered. Add one map or location figure and define the system boundary before moving into technical assumptions.}

\subsection{Design philosophy and requirements}
\instructionplaceholder{Required evidence}{State the design philosophy, performance objectives, relevant codes or standards, and the principal limit states that govern the concept. Define the design life, operational requirements, serviceability criteria, and the load combinations that must be checked. If a probabilistic philosophy is used, state why it is appropriate for this project.}

\subsection{Design conditions and load cases}
\instructionplaceholder{Required evidence}{List the governing design conditions that arise from the project function and site constraints, for example traffic actions, rail actions, inspection access, accidental events, or construction-stage requirements. Use a compact table to connect each load case to the design check it controls.}

\subsection{Metocean basis}
\instructionplaceholder{Required evidence}{Document the metocean data used for the project. Include water depth or bathymetry, wind climate, waves, currents, tides, and any other relevant hazards such as seismicity, landslides, ice, or long-term sea-level change. For each data source, state its origin, period of record, and why it is suitable for design use.}

\subsection{Environmental design values}
\instructionplaceholder{Required evidence}{Translate raw environmental data into design values. Show the Extreme Value Analysis or other method used to derive wave and environmental conditions, identify the governing return periods, and explain which combinations of wave, wind, and current control the design. Add geotechnical conditions here if they materially affect the structural response.}
